\subsection{Integration with respect to a measure defined by a density}

In the statements of this subsection, g denotes a positive numerical function, defined everywhere and locally \(\mu\) -integrable, \(\theta\) denotes a complex measure, and u a locally \(\theta\) -integrable complex function.

The remarks in No.~2 show that the results of §4 are applicable to the measure \(\nu = g \cdot \mu = \int g(t)\varepsilon_t d\mu(t)\) (even though the measure \(g(t)\varepsilon_t\) is not defined unless \(g(t) \neq +\infty\) ). We thus obtain the following statement:

PROPOSITION 3. --- For every numerical function \(f \geqslant 0\) defined on T,

\[
\int^ {\bullet} f d \nu = \int^ {\bullet} (f g) d \mu .\tag{5}
\]

This follows from Th. 1 of §4, No.~2.

COROLLARY 1. --- In order that a function f, with values in a Banach space or in \(\overline{R}\) , be locally negligible for the measure \(u \cdot \theta\) , it is necessary and sufficient that uf be locally negligible for \(\theta\) .

To say that f (resp. \(uf\) ) is locally negligible for \(u \cdot \theta\) (resp. for \(\theta\) ) is equivalent to saying that \(|f|\) (resp. \(|u| \cdot |f|\) ) is locally negligible for \(|u \cdot \theta|\) (resp. for \(|\theta|\) ). We are thus reduced, in view of Prop. 2 of No.~2, to the case that \(f, u, \theta\) are positive; the statement then follows at once from Prop. 3.

COROLLARY 2. --- Let \(u_1\) and \(u_2\) be two locally \(\theta\) -integrable complex functions. In order that \(u_1 \cdot \theta = u_2 \cdot \theta\) , it is necessary and sufficient that \(u_1\) and \(u_2\) be equal locally almost everywhere.

One is immediately reduced to showing that \(u \cdot \theta = 0\) implies that \(u(t) = 0\) locally almost everywhere; but \(u \cdot \theta = 0\) means that the function 1 is locally negligible for the measure \(u \cdot \theta\) . One then applies Corollary 1.

COROLLARY 3. --- Let u be a complex function that is locally integrable for the positive measure \(\mu\) . For \(u \cdot \mu\) to be a positive measure, it is necessary and sufficient that \(u(t) \geqslant 0\) locally almost everywhere.

For, \(u \cdot \mu\) is positive if and only if \(u \cdot \mu = |u \cdot \mu| = |u| \cdot \mu\) (Prop. 2), and this is equivalent to \(u = |u|\) locally almost everywhere (Cor. 2).

PROPOSITION 4. --- For a mapping f of T into a topological space G to be measurable for the measure \(u \cdot \theta\) , it is necessary and sufficient that the restriction of f, to the \(\theta\) -measurable set S of the t such that \(u(t) \neq 0\) , be \(\theta\) -measurable.

When u and \(\theta\) are positive, this follows at once from Prop. 3 of §4, No.~3. The result then extends to the case that u and \(\theta\) are complex thanks to Prop. 2.

COROLLARY. --- Let f be a function defined on T, with values in a Banach space F or in \(\overline{R}\) . For f to be \((u\cdot\theta)\) -measurable, it is necessary and sufficient that uf be \(\theta\) -measurable.

For, \(u\mathbf{f}\) is the extension by 0 of \((uf)|S\) to \(T\) .

THEOREM 1. --- Let f be a function defined on T, with values in a Banach space F or in \(\overline{R}\) . In order that f be essentially integrable for the measure \(\eta = u\cdot \theta\) , it is necessary and sufficient that uf be essentially \(\theta\) -integrable, in which case

\[
\int \mathbf {f} d \eta = \int (u \mathbf {f}) d \theta .\tag{6}
\]

Suppose moreover that \(u\) is continuous and that \(u(t) \neq 0\) for all \(t \in \mathbb{T}\) ; then \(\mathbf{f}\) is integrable for the measure \(\eta\) if and only if \(u\mathbf{f}\) is integrable for \(\theta\) .

The case that u and \(\theta\) are positive follows at once from Th. 2 of §4, No.~4. The first and the last assertion of the statement then follow immediately, because f is essentially integrable (resp. integrable) with respect to \(\eta = u \cdot \theta\) if and only if it is essentially integrable (resp. integrable) for \(|\eta| = |u| \cdot |\theta|\) . Finally, suppose that f is essentially integrable for \(\eta\) (hence for \(|\eta|\) ); we make use of the decomposition (2): f is essentially integrable for each of the measures \(\eta_{ij} = g_i \cdot \mu_j\) (i = 1, 2, 3, 4, j = 1, 2, 3, 4), because these measures are \(\leqslant |\eta|\) . We have

\[
\int \mathbf {f} d \eta_ {i j} = \int g _ {i} \mathbf {f} d \mu_ {j}.
\]

The formula (6) follows immediately from this.

COROLLARY. --- For the measure \(u \cdot \theta\) to be bounded, it is necessary and sufficient that u be essentially \(\theta\) -integrable.

Example. --- Let A be a subset of T; for \(\varphi_{A}\) to be locally \(\mu\) -integrable, it is necessary and sufficient that A be \(\mu\) -measurable. Assuming the condition to be fulfilled, set \(\nu = \varphi_{A} \cdot \mu\) ; for every numerical function \(f \geqslant 0\) defined on T, we then have

\[
\int^ {\bullet} f d \nu = \int^ {\bullet} f \varphi_ {\mathrm{A}} d \mu ,
\]

a quantity that is also denoted by \(\int_{\mathrm{A}}^{\bullet}f d\mu\) (cf.~Ch. IV, §5, No.~6). For a mapping \(g\) of T into a topological space G to be \(\nu\) -measurable, it is necessary and sufficient that the restriction of \(g\) to A be \(\mu\) -measurable. For a mapping \(\mathbf{f}\) of T into a Banach space F or into \(\overline{\mathbf{R}}\) to be essentially \(\nu\) -integrable, it is necessary and sufficient that \(\mathbf{f}\varphi_{\mathrm{A}}\) be essentially \(\mu\) -integrable, in which case

\[
\int \mathbf {f} d \nu = \int \mathbf {f} \varphi_ {\mathrm{A}} d \mu ,
\]

an expression that is also denoted \(\int_{A} f d\mu\) . Note that if two mappings of T into G (resp. F, \(\overline{R}\) ) coincide on A, then for one of them to be \(\nu\) -measurable (resp. essentially \(\nu\) -integrable), it is necessary and sufficient that the other be so. If now g is a mapping into G of an arbitrary subset \(B \supset A\) of T, one says that g is \(\mu\) -measurable on A if an arbitrary extension to T of the restriction of g to A is \(\nu\) -measurable, which amounts to saying that the restriction of g to A is \(\mu\) -measurable. A mapping f of B into a Banach space F, or into \(\overline{R}\) , is said to be essentially \(\mu\) -integrable on A if some extension \(\overline{f}\) to T of the restriction of f to A is essentially \(\nu\) -integrable; one then sets

\[
\int_ {\mathrm{A}} \mathbf {f} d \mu = \int_ {\mathrm{A}} \overline {{\mathbf {f}}} d \mu = \int \overline {{\mathbf {f}}} \varphi_ {\mathrm{A}} d \mu ,
\]

and one says that \(\int_{\mathrm{A}} \mathbf{f} d\mu\) is the integral of \(\mathbf{f}\) on \(\mathrm{A}\) (or extended to \(\mathrm{A}\) ). If \(f\) is a numerical function \(\geqslant 0\) defined on \(\mathrm{B} \supset \mathrm{A}\) , one defines similarly \(\int_{\mathrm{A}}^{*} f d\mu\) and \(\int_{\mathrm{A}}^{\bullet} f d\mu\) . Finally, a numerical function \(g\) defined on \(\mathrm{B} \supset \mathrm{A}\) is said to be locally \(\mu\) -integrable on \(\mathrm{A}\) if some extension \(\overline{g}\) to \(\mathrm{T}\) of the restriction of \(g\) to \(\mathrm{A}\) is locally \(\nu\) -integrable: this is equivalent to saying that, for every compact subset \(\mathrm{K}\) of \(\mathrm{T}\) , \(\overline{g}\varphi_{\mathrm{K} \cap \mathrm{A}}\) is \(\mu\) -integrable.

